\rowcolors{0}{}{}
\arrayrulecolor{black}
\renewcommand{\arraystretch}{1.3}
\setlength{\extrarowheight}{3pt}
\small
\sloppy
\begin{longtable}{|C{1cm}|L{6.5cm}|L{7cm}|C{1cm}|}
\caption{Backlog du sprint 4}\label{tab:backlog_sprint4}\\
\hline
\textbf{ID US} & \textbf{User Story} & \textbf{Description technique} & \textbf{Effort} \\
\hline
\endfirsthead

\hline
\textbf{ID US} & \textbf{User Story} & \textbf{Description technique} & \textbf{Effort} \\
\hline
\endhead

\hline
\endfoot

\hline
\endlastfoot

% Epic 13: Calendrier et planification
\multicolumn{4}{|l|}{\cellcolor{blue!20}\textbf{Epic 13 - Calendrier et planification}} \\
\hline

US66 & En tant qu'\textbf{utilisateur}, je souhaite visualiser les tickets et demandes de maintenance sur un calendrier mensuel afin d'avoir une vue globale des interventions. &
\begin{itemize}[nosep, left=0pt, label={$\bullet$}]
\item Mettre en place la page calendrier côté Next.js
\item S'appuyer sur une vue mois et une vue agenda/Gantt
\item Agréger les tickets, demandes et données de charge existantes
\item Colorer les événements selon leur type et leur statut
\item Permettre la navigation mois par mois et la consultation mobile
\end{itemize} & 3 \\
\hline

US67 & En tant qu'\textbf{administrateur}, je souhaite identifier visuellement les urgences SLA sur le calendrier afin de prioriser les interventions critiques. &
\begin{itemize}[nosep, left=0pt, label={$\bullet$}]
\item Calculer l'urgence à partir des échéances SLA des tickets
\item Mettre en évidence les tickets en retard ou proches de l'échéance
\item Utiliser une coloration visuelle et des badges d'alerte
\item Ajouter une légende explicative dans l'interface
\item Réutiliser les informations SLA déjà calculées côté backend
\end{itemize} & 2 \\
\hline

US68 & En tant que \textbf{chef de projet}, je souhaite accéder à une vue Gantt / Agenda afin d'analyser la charge de travail et la progression de chaque développeur. &
\begin{itemize}[nosep, left=0pt, label={$\bullet$}]
\item Intégrer la vue agenda/Gantt dans la page calendrier
\item Développer une visualisation de la charge par développeur
\item S'appuyer sur les endpoints existants de suivi de charge et de planning
\item Afficher la timeline des tickets assignés
\item Ajouter un filtre par développeur, projet et période
\end{itemize} & 3 \\
\hline

US69 & En tant qu'\textbf{utilisateur}, je souhaite accéder aux détails d'un ticket directement depuis le calendrier afin de consulter ou modifier l'intervention rapidement. &
\begin{itemize}[nosep, left=0pt, label={$\bullet$}]
\item Rendre les événements du calendrier cliquables
\item Ouvrir un panneau latéral ou rediriger vers la fiche ticket
\item Afficher le statut, la priorité, l'assignation et l'échéance
\item Permettre l'accès rapide à la vue complète du ticket
\end{itemize} & 2 \\
\hline

% Epic 10: Messagerie et collaboration
\multicolumn{4}{|l|}{\cellcolor{blue!20}\textbf{Epic 10 - Messagerie et collaboration}} \\
\hline

US70 & En tant qu'\textbf{utilisateur}, je souhaite accéder à une messagerie intégrée afin de communiquer avec mon équipe. &
\begin{itemize}[nosep, left=0pt, label={$\bullet$}]
\item Mettre en place les channels de discussion côté backend et frontend
\item Différencier les channels généraux, projets, tickets, demandes et privés
\item Utiliser les WebSockets pour la diffusion en temps réel
\item Créer l'interface de messagerie côté Next.js
\item Afficher les canaux, les membres et les messages non lus
\end{itemize} & 3 \\
\hline

US71 & En tant qu'\textbf{utilisateur}, je souhaite envoyer des messages dans un channel afin de communiquer. &
\begin{itemize}[nosep, left=0pt, label={$\bullet$}]
\item Développer l'endpoint \texttt{POST /api/messages/channel/:channelId}
\item Enregistrer les messages avec horodatage et auteur
\item Diffuser le message en temps réel via WebSocket
\item Afficher l'historique des messages du channel
\item Gérer les messages avec pièces jointes via un endpoint dédié
\end{itemize} & 2 \\
\hline

US72 & En tant qu'\textbf{utilisateur}, je souhaite répondre à un message afin de créer un fil de discussion. &
\begin{itemize}[nosep, left=0pt, label={$\bullet$}]
\item S'appuyer sur le champ \texttt{parentMessageId} dans le message
\item Regrouper les réponses en fil de discussion dans l'interface
\item Notifier les participants concernés dans le thread
\item Permettre le repli et le dépli des discussions
\end{itemize} & 2 \\
\hline

US73 & En tant qu'\textbf{utilisateur}, je souhaite mentionner des collègues afin d'attirer leur attention. &
\begin{itemize}[nosep, left=0pt, label={$\bullet$}]
\item Mettre en place l'autocomplétion des utilisateurs dans l'éditeur
\item Rechercher les utilisateurs pour la suggestion des mentions
\item Notifier les utilisateurs mentionnés via in-app et email
\item Mettre en évidence les messages contenant des mentions
\end{itemize} & 2 \\
\hline

US74 & En tant qu'\textbf{utilisateur}, je souhaite éditer mes messages afin de corriger des erreurs. &
\begin{itemize}[nosep, left=0pt, label={$\bullet$}]
\item Développer l'endpoint \texttt{PATCH /api/messages/:messageId}
\item Autoriser l'édition uniquement par l'auteur du message
\item Limiter l'édition à une fenêtre temporelle courte
\item Conserver l'historique des modifications
\item Afficher un indicateur de modification
\end{itemize} & 1 \\
\hline

US75 & En tant qu'\textbf{utilisateur}, je souhaite supprimer mes messages afin de retirer du contenu inapproprié. &
\begin{itemize}[nosep, left=0pt, label={$\bullet$}]
\item Développer l'endpoint \texttt{DELETE /api/messages/:messageId}
\item Confirmer l'action avant suppression
\item Conserver la suppression sous forme de soft delete
\item Préserver l'audit et l'historique si nécessaire
\end{itemize} & 1 \\
\hline

US76 & En tant qu'\textbf{utilisateur}, je souhaite épingler des messages importants afin de les retrouver facilement. &
\begin{itemize}[nosep, left=0pt, label={$\bullet$}]
\item Développer l'endpoint \texttt{POST /api/messages/:messageId/pin}
\item Afficher les messages épinglés en haut du channel
\item Ajouter une icône d'épingle sur le message
\item Restreindre cette action selon le rôle dans le channel
\end{itemize} & 1 \\
\hline

US77 & En tant que \textbf{chef de projet}, je souhaite voir qui est en ligne dans mon équipe afin de savoir qui je peux solliciter. &
\begin{itemize}[nosep, left=0pt, label={$\bullet$}]
\item Utiliser le suivi de dernière connexion et les événements temps réel
\item Afficher le statut de présence dans les listes d'utilisateurs
\item Mettre en évidence les membres récemment actifs
\item Présenter les dernières activités dans l'interface
\end{itemize} & 2 \\
\hline

US78 & En tant qu'\textbf{utilisateur}, je souhaite voir qui est en ligne afin de savoir qui est disponible. &
\begin{itemize}[nosep, left=0pt, label={$\bullet$}]
\item Afficher les statuts de présence dans les composants de messagerie
\item Synchroniser l'état via les événements temps réel
\item Montrer la dernière activité ou la dernière connexion
\item Adapter l'affichage selon le rôle et le contexte
\end{itemize} & 2 \\
\hline

% Epic 11: Gestion des documents
\multicolumn{4}{|l|}{\cellcolor{blue!20}\textbf{Epic 11 - Gestion des documents}} \\
\hline

US79 & En tant qu'\textbf{utilisateur}, je souhaite uploader des documents afin de centraliser les fichiers. &
\begin{itemize}[nosep, left=0pt, label={$\bullet$}]
\item Développer l'endpoint \texttt{POST /api/documents/upload}
\item Autoriser jusqu'à 5 fichiers par envoi
\item Limiter la taille à 10 MB par fichier
\item Stocker les fichiers localement dans le répertoire documents
\item Associer les documents à un ticket ou à un contrat
\end{itemize} & 3 \\
\hline

US80 & En tant qu'\textbf{utilisateur}, je souhaite télécharger des documents afin de les consulter hors ligne. &
\begin{itemize}[nosep, left=0pt, label={$\bullet$}]
\item Prévoir le téléchargement à partir des fichiers stockés
\item Tracer la consultation des documents pour la traçabilité
\item Rattacher le document à son ticket ou contrat d'origine
\item Préparer l'ouverture dans le navigateur ou le téléchargement direct
\end{itemize} & 2 \\
\hline

US81 & En tant qu'\textbf{utilisateur}, je souhaite gérer les versions des documents afin d'assurer la traçabilité. &
\begin{itemize}[nosep, left=0pt, label={$\bullet$}]
\item Développer l'endpoint \texttt{GET /api/documents/:id/versions}
\item Développer l'endpoint \texttt{PATCH /api/documents/:id/replace}
\item Conserver les anciennes versions lors du remplacement
\item Exiger un commentaire lors du remplacement
\item Permettre l'accès à l'historique des versions
\end{itemize} & 2 \\
\hline

US82 & En tant qu'\textbf{utilisateur}, je souhaite consulter l'historique des consultations afin de voir qui a accédé au document. &
\begin{itemize}[nosep, left=0pt, label={$\bullet$}]
\item Développer l'endpoint \texttt{GET /api/documents/:id/consultations}
\item Afficher les consultations avec date, heure et utilisateur
\item Prévoir un filtre par période
\item Permettre l'exploitation de l'historique à des fins de traçabilité
\end{itemize} & 1 \\
\hline

US83 & En tant qu'\textbf{utilisateur}, je souhaite archiver des documents afin de nettoyer la vue active. &
\begin{itemize}[nosep, left=0pt, label={$\bullet$}]
\item Développer l'endpoint \texttt{PATCH /api/documents/:id/archive}
\item Développer l'endpoint \texttt{PATCH /api/documents/:id/restore}
\item Exclure les documents archivés de la vue courante
\item Réserver l'archivage à l'administrateur
\end{itemize} & 1 \\
\hline

US84 & En tant que \textbf{chef de projet}, je souhaite générer automatiquement un PV à partir d'un ticket afin de documenter les interventions. &
\begin{itemize}[nosep, left=0pt, label={$\bullet$}]
\item Générer un PDF à partir des informations du ticket
\item Inclure les actions réalisées, les intervenants et le temps passé
\item Enregistrer le PV comme document lié au ticket
\item Prévoir une exportation ou une signature selon le besoin métier
\end{itemize} & 3 \\
\hline

US85 & En tant que \textbf{client}, je souhaite télécharger les PV de mes tickets afin de conserver une trace des interventions. &
\begin{itemize}[nosep, left=0pt, label={$\bullet$}]
\item Exposer les PV liés aux tickets du client dans son espace
\item Permettre le téléchargement des PV générés
\item Notifier le client lors de la génération d'un nouveau PV
\item Rattacher le PV au contrat ou au ticket concerné
\end{itemize} & 1 \\
\hline

US86 & En tant que \textbf{développeur ou chef de projet}, je souhaite consulter la base de connaissances afin de trouver des solutions existantes. &
\begin{itemize}[nosep, left=0pt, label={$\bullet$}]
\item Mettre en place la recherche dans la base de connaissances
\item Filtrer par mots-clés, catégorie et tags
\item Afficher le titre, l'extrait et le score de pertinence
\item Ouvrir l'article complet en un clic
\end{itemize} & 2 \\
\hline

US87 & En tant qu'\textbf{administrateur}, je souhaite créer des articles de la base de connaissances afin d'enrichir la documentation. &
\begin{itemize}[nosep, left=0pt, label={$\bullet$}]
\item Développer l'endpoint de création d'articles KB
\item Utiliser un éditeur Markdown pour la rédaction
\item Ajouter des tags, une catégorie et un auteur
\item Permettre l'archivage des articles obsolètes
\end{itemize} & 2 \\
\hline

US88 & En tant qu'\textbf{utilisateur}, je souhaite créer une nouvelle version d'un document afin de maintenir l'historique. &
\begin{itemize}[nosep, left=0pt, label={$\bullet$}]
\item Réutiliser le mécanisme de remplacement de document
\item Conserver toutes les versions précédentes
\item Numéroter implicitement les versions dans l'historique
\item Exiger un commentaire pour documenter la nouvelle version
\end{itemize} & 2 \\
\hline

US89 & En tant qu'\textbf{utilisateur}, je souhaite comparer deux versions d'un document afin de voir les différences. &
\begin{itemize}[nosep, left=0pt, label={$\bullet$}]
\item Prévoir un endpoint ou un traitement de comparaison de versions
\item Afficher les deux versions côte à côte
\item Mettre en évidence les différences textuelles
\item Faciliter la revue des changements avant validation
\end{itemize} & 2 \\
\hline

US90 & En tant qu'\textbf{administrateur}, je souhaite voir qui a consulté un document afin de suivre sa diffusion. &
\begin{itemize}[nosep, left=0pt, label={$\bullet$}]
\item Exploiter l'historique des consultations de documents
\item Afficher les utilisateurs, dates et périodes de consultation
\item Prévoir une vue statistique pour l'administration
\item Permettre l'export des données si nécessaire
\end{itemize} & 2 \\
\hline

% Epic 12: Notifications
\multicolumn{4}{|l|}{\cellcolor{blue!20}\textbf{Epic 12 - Notifications et alertes}} \\
\hline

US91 & En tant qu'\textbf{utilisateur}, je souhaite recevoir des notifications in-app afin d'être informé en temps réel. &
\begin{itemize}[nosep, left=0pt, label={$\bullet$}]
\item Développer l'API de notification in-app
\item Afficher la liste paginée des notifications
\item Gérer le compteur de notifications non lues
\item Permettre le marquage lu et lu partout
\item Rendre les notifications cliquables
\end{itemize} & 2 \\
\hline

US92 & En tant qu'\textbf{utilisateur}, je souhaite recevoir des notifications par email afin d'être informé hors connexion. &
\begin{itemize}[nosep, left=0pt, label={$\bullet$}]
\item Configurer l'envoi d'emails transactionnels
\item S'appuyer sur une file de traitement asynchrone
\item Envoyer des emails pour les événements critiques
\item Prévoir la configuration de préférences de notification
\end{itemize} & 2 \\
\hline

US93 & En tant que \textbf{membre technique}, je souhaite recevoir des notifications Slack afin d'être alerté sur mon outil de travail. &
\begin{itemize}[nosep, left=0pt, label={$\bullet$}]
\item Intégrer un service de notification Slack
\item Envoyer des alertes pour les événements critiques
\item Prévoir la configuration des canaux selon le besoin métier
\item Compléter les notifications in-app et email
\end{itemize} & 2 \\
\hline

US94 & En tant qu'\textbf{utilisateur}, je souhaite recevoir des notifications push navigateur afin d'être alerté même hors onglet. &
\begin{itemize}[nosep, left=0pt, label={$\bullet$}]
\item Mettre en place les notifications push navigateur
\item Demander et gérer l'autorisation utilisateur
\item Envoyer des push pour les événements critiques
\item Ouvrir directement l'application lors du clic
\end{itemize} & 2 \\
\hline

US95 & En tant qu'\textbf{utilisateur}, je souhaite consulter l'historique de mes notifications afin de retrouver une information. &
\begin{itemize}[nosep, left=0pt, label={$\bullet$}]
\item Développer la liste paginée des notifications
\item Ajouter un filtre par statut lu ou non lu
\item Prévoir un filtre temporel dans l'interface
\item Permettre la recherche dans l'historique
\end{itemize} & 1 \\
\hline

% Epic 13: Dashboards et rapports
\multicolumn{4}{|l|}{\cellcolor{blue!20}\textbf{Epic 13 - Dashboards et rapports}} \\
\hline

US96 & En tant qu'\textbf{administrateur}, je souhaite consulter un dashboard global afin de piloter le système. &
\begin{itemize}[nosep, left=0pt, label={$\bullet$}]
\item Développer l'endpoint \texttt{GET /api/dashboard/stats}
\item Afficher les KPI globaux du système
\item Suivre les tickets, les contrats, les clients et la charge
\item Ajouter les alertes SLA et les tendances récentes
\end{itemize} & 3 \\
\hline

US97 & En tant que \textbf{chef de projet}, je souhaite consulter un dashboard d'équipe afin de piloter l'activité. &
\begin{itemize}[nosep, left=0pt, label={$\bullet$}]
\item Exploiter les données de tickets et de charge développeur
\item Afficher la répartition de la charge par membre
\item Suivre les retards, les urgences et les tickets bloqués
\item Prévoir une vue orientée pilotage d'équipe
\end{itemize} & 2 \\
\hline

US98 & En tant que \textbf{développeur}, je souhaite consulter mon dashboard personnel afin d'organiser ma journée. &
\begin{itemize}[nosep, left=0pt, label={$\bullet$}]
\item Développer l'endpoint \texttt{GET /api/tickets/dashboard/employee}
\item Afficher les tickets assignés avec priorité et échéance
\item Lister les urgences du jour
\item Suivre le temps de travail journalier
\item Identifier les échéances proches et les tickets en cours
\end{itemize} & 2 \\
\hline

US99 & En tant que \textbf{client}, je souhaite consulter mon dashboard afin de suivre mes projets. &
\begin{itemize}[nosep, left=0pt, label={$\bullet$}]
\item Développer l'endpoint \texttt{GET /api/dashboard/client}
\item Afficher les contrats actifs et leur état
\item Suivre les demandes et tickets ouverts
\item Présenter les documents et informations utiles au client
\end{itemize} & 2 \\
\hline

US100 & En tant qu'\textbf{administrateur}, je souhaite générer un rapport d'activité mensuel afin de présenter les statistiques. &
\begin{itemize}[nosep, left=0pt, label={$\bullet$}]
\item Développer l'endpoint \texttt{GET /api/dashboard/rapport-mensuel}
\item Calculer les tickets traités, le temps passé et le respect SLA
\item Prévoir l'export PDF et Excel côté interface
\item Ajouter des graphiques pour la présentation des résultats
\end{itemize} & 3 \\
\hline

US101 & En tant qu'\textbf{administrateur}, je souhaite générer des rapports de conformité SLA afin de justifier les performances. &
\begin{itemize}[nosep, left=0pt, label={$\bullet$}]
\item Exploiter les indicateurs SLA du dashboard admin
\item Calculer le taux de conformité par période ou par contrat
\item Identifier les tickets en dépassement
\item Prévoir un export de rapport pour la justification métier
\end{itemize} & 2 \\
\hline

US102 & En tant que \textbf{chef de projet}, je souhaite générer des rapports d'activité afin de présenter l'avancement. &
\begin{itemize}[nosep, left=0pt, label={$\bullet$}]
\item Agréger les données de tickets et de charge d'équipe
\item Afficher l'évolution de l'activité par période
\item Comparer les tickets créés, traités et en retard
\item Prévoir un export PDF ou Excel selon le besoin
\end{itemize} & 2 \\
\hline

\end{longtable}